\documentclass[12pt, addpoints]{exam}
\usepackage[utf8]{inputenc}
\usepackage[T1]{fontenc}
\usepackage[brazilian]{babel}
\usepackage{amsmath, amssymb}
\usepackage{graphicx}
\usepackage{tikz}
\usepackage{pgfplots}
\pgfplotsset{compat=1.18}
\usepackage{booktabs}
\usepackage{xcolor}
\usepackage{mdframed}
\usepackage[top=1.5cm, bottom=1.5cm, left=2cm, right=2cm, headheight=14pt, headsep=8pt, footskip=10pt, includehead, includefoot]{geometry}

\marksnotpoints
\pointsinrightmargin
\marginpointname{ pts}

% Configuracao de Cabecalho e Rodape
\pagestyle{headandfoot}
\runningheader{Álgebra Linear}{}{Página \thepage\ de \numpages}
\footer{}{}{}



\begin{document}

% -----------------------------------------------------------
% Cabecalho
% -----------------------------------------------------------
\begin{center}
    \textsc{\textbf{Universidade Federal do Pará} \\
    CAMTUC -- FEE: Álgebra Linear} -- Prof.\ Dr.\ Raphael Teixeira \\
    \vspace{0.1cm}
    Avaliação 2: Determinantes e Espaços Vetoriais \hspace{0.6cm}-\hspace{0.6cm} 22/06/2026.
\end{center}
\vspace{0.3cm}
\noindent
\textbf{Aluno(a):} \rule{8cm}{0.1pt} \quad \textbf{Matrícula:} \rule{3.5cm}{0.1pt}
\vspace{0.0cm}
% -----------------------------------------------------------
\begin{questions}

% ============================================================
%  QUESTAO 1 -- Determinantes / Regra de Cramer
% ============================================================
\question
Considere os dois sistemas lineares abaixo:
\[
  \text{(I)}\quad
  \begin{cases}
    \phantom{2}x - \phantom{2}y + 2z = 7 \\[2pt]
    2x + \phantom{2}y - \phantom{2}z = 2 \\[2pt]
    \phantom{2}x + 2y + \phantom{2}z = 7
  \end{cases}
  \qquad\qquad
  \text{(II)}\quad
  \begin{cases}
    \phantom{2}x + \phantom{2}y + \phantom{2}z = 4 \\[2pt]
    2x - \phantom{2}y + 3z = 3 \\[2pt]
    3x \phantom{{}+y{}} + 4z = 5
  \end{cases}
\]

\textbf{(a, 1.5 pts)} Verifique, por meio de determinante, qual dos dois possui solução única. Justifique. \textbf{(b, 1.5 pt)} Para o sistema que admite solução única, utilize a \textbf{regra de Cramer} para determinar o valor de $y$.

% ============================================================
%  QUESTAO 2 -- Independencia Linear, Base e Dimensao (4.4 e 4.5)
% ============================================================


\question
Considere os vetores $v_1=(3,\,1,\,-1)$, $v_2=(2,\,1,\,0)$ e $v_3=(0,\,1,\,2)$, em $\mathbb{R}^3$.

\textbf{(a, 1.5 pts)} Verifique se o conjunto $\{v_1, v_2, v_3\}$ é linearmente dependente. Em caso afirmativo, exiba uma relação de dependência linear entre eles.

\textbf{(b, 1.5 pts)} Mostre que $u=(-1,\,-1,\,-1)$ pertence ao subespaço gerado por $\{v_1, v_2, v_3\}$, então escreva $u$ como combinação da base deste subespaço.

% ============================================================
%  QUESTAO 3 -- Posto, Espaco Coluna, Espaco Nulo e Solucao Geral (4.6)
% ============================================================
\question
Considere o sistema linear $A\mathbf{x} = \mathbf{b}$ dado por
\[
  \begin{cases}
    \phantom{2}x_1 \phantom{{}+x_2{}} + 2x_3 - \phantom{2}x_4 = 3\\[3pt]
    \phantom{2}x_1 + \phantom{2}x_2 + \phantom{2}x_3 + 2x_4 = 5\\[3pt]
    2x_1 + \phantom{2}x_2 + 3x_3 + \phantom{2}x_4 = 8
  \end{cases}
\]

\textbf{(a, 1.0 pt)} Determine o \textbf{posto} e a \textbf{nulidade} de $A$.

\textbf{(b, 0.5 pt)} Indique, entre os vetores-coluna \emph{originais} de $A$, quais formam uma \textbf{base do espaço coluna} de $A$.

\textbf{(c, 1.5 pt)} Determine uma \textbf{solução particular} $\mathbf{x}_p$ do sistema $A\mathbf{x} = \mathbf{b}$ e escreva a \textbf{solução geral} do sistema na forma $\mathbf{x} = \mathbf{x}_p + \mathbf{x}_h$.

\textbf{(d, 1.0 pt)} Determine uma \textbf{base do espaço nulo} de $A$.

\end{questions}
\end{document}